% =====================================================================
%  Conferenze, GAMM Roma (Checkpoint Festival, 3a edizione) — Panel 4
%  Preservare il videogioco: archivi, musei e mercato in Europa
%  Fonte: trascrizioni automatiche "Track 4" e "Track 5", che sono due
%  parti dello stesso panel, in inglese e in parte in italiano, con gli
%  ospiti collegati da remoto. La videochiamata cade a metà
%  dell'intervento del rappresentante di GOG: la Track 4 termina con
%  ripetizioni meccaniche prive di senso, la Track 5 riprende a frase
%  iniziata dopo la riconnessione.
%  Etichette con prefisso "c1-4:". Nessun pacchetto aggiuntivo.
% =====================================================================

\chapter[Preservare il videogioco]{Preservare il videogioco: archivi, musei e
mercato in Europa}
\label{cap:c1-4-preservazione}

Il quarto incontro al GAMM è un panel, in gran parte \textbf{in inglese}, dedicato
alla \textbf{preservazione del videogioco}. È organizzato insieme a
\textbf{EFGAMP}\footnote{Trascritto ``EFGAP'', ``FCAMP'', ``Afghan'',
``FGAMP'', ``Ed Gamper''. EFGAMP è la European Federation of Game Archives,
Museums and Preservation Projects.}, la federazione europea di archivi, musei e
progetti di preservazione del videogioco, di cui il GAMM fa parte (``anche se non
siamo tra i fondatori''). La moderatrice è la curatrice del museo\footnote{Gli
ospiti la chiamano ``Katharina''; nel capitolo~\ref{cap:c1-2-videogioco-storico}
la curatrice si presenta come ``Caterina''.}. A causa di uno
\textbf{sciopero}, il programma della giornata è stato cambiato e il panel è
trasmesso in diretta sulla pagina Instagram del museo; tutti e tre gli ospiti
sono \textbf{collegati da remoto}, e a metà la videochiamata cade, per poi
riprendere\footnote{Le due trascrizioni ``Track 4'' e ``Track 5'' sono le due
metà del panel. La caduta della chiamata produce nella prima una serie di frasi
ripetute senza senso, che non si riportano.}.

Gli ospiti rappresentano, dice la moderatrice, \textbf{tre approcci} diversi alla
preservazione: uno più \textbf{storico}, uno più \textbf{pubblico}, uno più
\textbf{commerciale}.
\begin{description}
  \item[Andreas Lange] dirige EFGAMP ed è stato il fondatore del
        \textbf{Computerspielemuseum di Berlino}\footnote{Trascritto
        ``Computer Studio di Berlino''.}, il primo museo del videogioco in
        Europa. Era presente all'inaugurazione del GAMM, due anni prima.
  \item[Willem] è ricercatore\footnote{Si presenta come ``game scientist''
        (forse ``game researcher''); il cognome è trascritto ``Hilworth'' e
        ``Pelhorst'' e non è verificabile.} e responsabile dei giochi e dei
        media online presso il \textbf{Netherlands Institute for Sound \&
        Vision} (l'istituto audiovisivo dei Paesi Bassi, a Hilversum), membro del
        board di EFGAMP.
  \item[Piotr] è membro del board di EFGAMP e \emph{senior communication and
        brand advocacy manager} di \textbf{GOG}\footnote{Trascritto ``DOG'',
        ``GLG'', ``Gio G'', ``GUG''. GOG è il negozio digitale di giochi senza DRM
        del gruppo CD Projekt. Il cognome di Piotr è trascritto ``Knibner'' e
        non è verificabile.}; la moderatrice lo ha conosciuto un mese prima a un
        evento di settore a Bologna.
\end{description}

Struttura del capitolo:
\begin{enumerate}
  \item le sfide europee secondo EFGAMP (sezione~\ref{sec:c1-4-sfide});
  \item il progetto \emph{Games2Cloud} e l'ontologia delle collezioni
        (sezione~\ref{sec:c1-4-cloud});
  \item segnali dal mondo e dall'Europa (sezione~\ref{sec:c1-4-segnali});
  \item l'archivio pubblico olandese e il progetto per gli sviluppatori
        (sezione~\ref{sec:c1-4-olanda});
  \item il punto di vista commerciale: GOG (sezione~\ref{sec:c1-4-gog});
  \item perché preservare conviene anche all'industria
        (sezione~\ref{sec:c1-4-business});
  \item la domanda sul futuro del supporto fisico
        (sezione~\ref{sec:c1-4-fisico});
  \item i premi (sezione~\ref{sec:c1-4-premi}).
\end{enumerate}

% ---------------------------------------------------------------------
\section{Le sfide della preservazione in Europa}
\label{sec:c1-4-sfide}

Prima domanda ad Andreas Lange: quali sono oggi le principali sfide europee?

La situazione, risponde, ``è ottima'' sotto un aspetto: l'Europa ha qualcosa di
\textbf{unico a livello globale}, una federazione con membri da tutti i settori
culturali che si occupano di collezionare e preservare giochi. EFGAMP ha
\textbf{30 membri}, ``dalle iniziative dal basso alle biblioteche nazionali e
tutto ciò che sta in mezzo''. Avere una struttura organizzata è un risultato
importante; altrove nel mondo stanno nascendo iniziative simili, ed è sicuro che
tra qualche anno EFGAMP sarà ``solo una tra le altre''.

\subsection{Comunità e istituzioni}

La \textbf{comunità dei giocatori} è stata finora un attore cruciale della
preservazione, e lo resterà: basti pensare agli \textbf{emulatori}, programmati e
mantenuti da appassionati. Ma le generazioni cambiano: chi ha un legame emotivo
con i giochi più vecchi tra cinquant'anni forse non ci sarà più. Per questo
diventano importanti le \textbf{istituzioni}, che sono più \textbf{sostenibili}
e più \textbf{sistematiche}.

\subsection{Due problemi}

\begin{enumerate}
  \item \textbf{Leggi sul diritto d'autore frammentate.} Esiste un quadro
        europeo, ma è recepito in modo diverso dalle leggi nazionali degli stati
        membri. Questo rende più complesso preservare i giochi e soprattutto
        \textbf{renderli accessibili}, che è lo scopo della preservazione.
  \item \textbf{Il passaggio dal supporto fisico al gioco su server.} Gli
        scaffali delle istituzioni sono pieni di scatole, floppy disk, CD-ROM:
        una copia fisica \textbf{la possiedi} e con essa puoi fare molte cose.
        Con un gioco che gira solo su un server, e di cui hai solo un accesso,
        no. È un cambiamento che rende la preservazione ``molto, molto più
        complicata''.
\end{enumerate}

La campagna \textbf{Stop Killing Games} mostra che la consapevolezza del problema
sta crescendo: ha raggiunto la soglia di un'\textbf{iniziativa dei cittadini
europei}, e il Parlamento europeo dovrà discuterne, non tanto in prospettiva
culturale quanto di \textbf{diritti dei consumatori}. Ma è la stessa cosa: se i
consumatori non possono più giocare, anche le istituzioni hanno un grosso
problema quando il detentore dei diritti ``spegne'' il gioco.

% ---------------------------------------------------------------------
\section{Games2Cloud e un'ontologia per le collezioni}
\label{sec:c1-4-cloud}

Una federazione serve a fare ciò che i singoli non possono: raccogliere i
bisogni e costruire una \textbf{strategia condivisa}. Lange condivide lo schermo
per mostrare un progetto appena finanziato a EFGAMP, \textbf{Games2Cloud}.

Il progetto si collega a una grande iniziativa dell'Unione europea, il
\textbf{Cultural Heritage Cloud}\footnote{Il nome ufficiale è European
Collaborative Cloud for Cultural Heritage (ECCCH).}, che diventerà una risorsa
per le istituzioni della memoria di tutti i settori:
\begin{itemize}
  \item l'UE investe \textbf{oltre 102 milioni di euro}, dal 2023 al 2029;
  \item il cloud è costruito attraverso molti progetti \textbf{Horizon}
        (ventidue in elenco), che forniscono funzionalità, coordinati da un
        progetto apposito\footnote{Trascritto ``ECOS''; si tratta verosimilmente
        del progetto di coordinamento ECHOES.};
  \item lo scopo è riunire comunità diverse, fornire loro strumenti e
        \textbf{collegare} dati e collezioni per offrirli a contesti più ampi.
\end{itemize}
Con Games2Cloud anche il settore del \textbf{patrimonio videoludico} è collegato
a questo cloud, cosa che nessun attore singolo avrebbe potuto ottenere. Il
pubblico di riferimento comprende anche le \textbf{comunità di pratica}: non
serve essere un museo, una biblioteca nazionale o un'istituzione GLAM (gallerie,
biblioteche, archivi, musei). La speranza è che, quando il cloud sarà pronto nel
2029, anche molte comunità di \emph{retrogaming} possano beneficiarne.

\subsection{L'ontologia}

Con i fondi europei EFGAMP sta prototipando un'\textbf{ontologia per le
collezioni di videogiochi}. Un'ontologia collega tra loro insiemi di dati
diversi e permette di cercare in molti database \textbf{con una sola
ricerca}: è il principio del web \textbf{semantico} e dei \emph{linked open
data}. Per le collezioni di videogiochi un'ontologia del genere \textbf{non esiste
ancora}. È un tassello importante della strategia: rendere i risultati delle
collezioni rilevanti in un contesto sociale più ampio.

% ---------------------------------------------------------------------
\section{Segnali dal mondo e dall'Europa}
\label{sec:c1-4-segnali}

C'è qualche progetto, anche fuori dall'Europa, da tenere d'occhio?

\paragraph{Il padiglione cinese alla Biennale di Venezia.} Proprio perché il
festival è in Italia, Lange segnala che il padiglione della Cina alla
\textbf{Biennale di Venezia} espone uno dei videogiochi cinesi più famosi,
\emph{Black Myth: Wukong}\footnote{Trascritto ``Black News Wukong'' e, più
avanti, ``Black Lives Matter'': in entrambi i casi si tratta di \emph{Black Myth:
Wukong}.}. Non gli risulta che un videogioco sia mai stato esposto prima in un
padiglione nazionale. L'anno precedente l'Art Museum della China Academy of Art,
un grande museo d'arte moderna, gli aveva dedicato un'intera mostra. È il segno
che in Cina l'immagine del videogioco presso l'\textbf{establishment culturale}
è ``molto, molto positiva''.

\paragraph{Uno studio pilota europeo.} Si è appena chiuso il bando per uno
\textbf{studio pilota sulla salvaguardia del patrimonio europeo nato digitale}
(\emph{born-digital heritage}), che diventerà la \textbf{roadmap} dei futuri
programmi europei. Il bando individua tre settori concreti: \textbf{media art},
\textbf{archiviazione del web} e \textbf{videogiochi}. Lange, ``ottimista da
trent'anni'', vede il tempo lavorare a favore della preservazione, pur con molto
lavoro ancora da fare.

% ---------------------------------------------------------------------
\section{Un archivio pubblico: il caso olandese}
\label{sec:c1-4-olanda}

La moderatrice osserva che EFGAMP unisce istituzioni con scopi diversi: un ente
pubblico come quello di Willem e un'azienda come GOG, ``importante partner''
nonostante le regole d'ingresso molto rigide.

\subsection{Che cosa e come si archivia}

Il Netherlands Institute for Sound \& Vision archivia videogiochi
\textbf{dal 2017--2018}, perché li considera \textbf{patrimonio culturale
digitale}, in particolare quelli \textbf{sviluppati nei Paesi Bassi}: ``non
collezioniamo \emph{The Legend of Zelda}, ma un gioco di corse molto famoso uscito
solo nei Paesi Bassi''\footnote{Trascritto ``A2 Racer'', una serie olandese di
giochi di corse ambientati sull'autostrada A2. Willem aggiunge che dallo stesso
motore lo studio ricavò anche un gioco di corse in motoscafo sui canali di
Venezia, ``perché avevano un solo motore e cercavano di sfruttarlo il più
possibile''.}.

I principi:
\begin{itemize}
  \item preservazione \textbf{digitale}, non fisica: i giochi sono \emph{born
        digital}, ``in fondo sono zeri e uni'' (``sentimentalmente penso che
        siano molto di più'', ammette, ma ai fini della preservazione si
        ragiona così);
  \item si chiede sempre il \textbf{permesso ai detentori dei diritti}: l'istituto
        dipende dalla comprensione di chi possiede i diritti, e conserva i giochi
        disponibili \textbf{all'interno dell'istituto} per scopi di
        \textbf{ricerca e riuso};
  \item l'istituto è prima di tutto un \textbf{archivio audiovisivo}: il suo
        compito principale è archiviare le emittenti radiotelevisive; i
        videogiochi si sono aggiunti perché ritenuti altrettanto importanti.
\end{itemize}

\subsection{Il problema: raggiungere gli sviluppatori}

Il mondo culturale può dire che i giochi sono patrimonio da preservare, ma
\textbf{far attecchire l'idea tra gli sviluppatori} è difficile: anche negli
studi, piccoli o grandi, spesso non si considera che cosa significhi
preservazione. E gli sviluppatori sono sempre impegnati sul progetto successivo:
se l'archivio arriva \textbf{dopo}, c'è una grande probabilità che conoscenze e
accessi ai materiali siano \textbf{perduti}.

Di qui un progetto lanciato da poche settimane, una pagina web chiamata
\textbf{Control Save}\footnote{Trascritto ``CTRLSA'', ``Control Safe'',
``ControlSave'', ``Cultural Save''; il gioco di parole è con la scorciatoia
\textsc{ctrl+s}. La grafia esatta non è verificabile.}, in olandese e in inglese:
una breve panoramica di \textbf{linee guida e punti di partenza} per gli
sviluppatori su che cosa considerare per preservare il proprio gioco. Distilla
le regole usate dall'istituto in una decina d'anni, così che gli sviluppatori
possano applicarle \textbf{senza dover contattare l'archivio}, e lavorare
``preventivamente''. Contiene molti \textbf{video esplicativi}.

\subsection{Il video sugli asset}

Willem mostra il video dedicato agli \textbf{asset di gioco}. Il contenuto:
\begin{itemize}
  \item le collezioni storiche sono di due tipi. Da un lato quelle che
        conservano \textbf{giochi funzionanti e software}, spesso privati o su
        macchine dedicate; dall'altro gli \textbf{istituti culturali}, che
        vogliono anche i \textbf{materiali di contesto};
  \item questi ultimi vogliono documentare la \textbf{storia dei media}: chi
        visita l'archivio vuole sapere come i giochi sono stati fatti, da chi,
        quando e in quali condizioni. Le domande possono essere ampie (``che tipo
        di giochi si facevano in quel periodo?'') o molto specifiche (``a che
        punto era il \emph{ray tracing} nello sviluppo di giochi nel 1999?'');
  \item perciò scegliere che cosa donare è complesso: tutte le iterazioni della
        documentazione o solo il file finale? Solo la \emph{key art} o il file di
        ogni tubo di scarico 3D mai modellato? Il codice, o solo demo, prototipi e
        build finali?
  \item non esiste \textbf{una regola d'oro}: ogni studio è diverso, e gli
        archivisti sono felici di esaminare con gli sviluppatori ciò che offrono
        e suggerire che cosa includere o escludere.
\end{itemize}
Molti sviluppatori, conclude Willem, non considerano archivi e preservazione
\textbf{parte del proprio lavoro}; il progetto vuole avviare questa idea.

\subsection{Perché in inglese}

Perché un progetto olandese anche in inglese? Nei Paesi Bassi lavora un
\textbf{grande gruppo internazionale} di sviluppatori, tra grandi studi e molto
\textbf{co-sviluppo}; e il messaggio deve arrivare lontano, perché molte
istituzioni in Europa e nel mondo stanno facendo pionieristicamente lo stesso
lavoro, ognuna a modo suo. Control Save non vuole essere ``la torre d'avorio''
che detta regole d'oro, ma \textbf{un punto di partenza per il dialogo}: chi
lavora in Italia può portare quelle idee al GAMM. La moderatrice aggiunge che
EFGAMP sta lavorando a qualcosa di simile con il contributo di varie
istituzioni europee.

\subsection{Le sfide tecniche per un umanista}

Willem ha un background in \textbf{media, giornalismo e ricerca}, non è un
programmatore (``probabilmente è una bestemmia in chiesa''). Le sfide tecniche
sono:
\begin{itemize}
  \item \textbf{digitalizzare i supporti fisici} (dischi, cassette);
  \item \textbf{preservare gli zeri e uni}: molti giochi non sono mai stati
        pensati per durare quanto stanno già durando;
  \item \textbf{catturare il contesto}, spesso trascurato: i giochi raccontano
        una storia importante, ma anche il periodo in cui sono stati sviluppati
        è essenziale per capire la storia più ampia.
\end{itemize}
La vera fatica è \textbf{chiedere aiuto} e trovare il tempo per farlo, sia agli
sviluppatori sia ai colleghi delle altre istituzioni (e, per le questioni
tecniche, a realtà come GOG).

% ---------------------------------------------------------------------
\section{Il punto di vista commerciale: GOG}
\label{sec:c1-4-gog}

Dopo la caduta della chiamata, la registrazione riprende mentre Piotr spiega che
le regole a cui deve sottostare un'azienda commerciale sono \textbf{più
stringenti} di quelle delle istituzioni GLAM: per GOG non esiste la nozione di
\emph{abandonware}, bisogna \textbf{sempre} trovare il proprietario della
proprietà intellettuale e verificare tutte le licenze del gioco.

\subsection{Perché i giochi scompaiono}

Soprattutto per ragioni legali:
\begin{itemize}
  \item un \textbf{accordo morto} tra le parti;
  \item diritti \textbf{divisi} tra più soggetti;
  \item un accordo \textbf{perduto}: dopo vent'anni o trenta capita che
        l'archivio sia andato bruciato e non ci sia una copia digitale; non puoi
        dimostrare che il gioco è tuo;
  \item le \textbf{licenze} contenute nel gioco: in titoli come \emph{FIFA} o
        \emph{Need for Speed}\footnote{Trascritto ``Infor Speeds''.} ci sono
        musica, automobili, sportivi e attori su licenza; ma anche
        \emph{middleware}, codec e altri componenti. Durante lo sviluppo si dice
        ``licenziamolo per dieci anni'', perché nessuno pensa a che cosa
        succederà tra venti, trenta o cinquant'anni.
\end{itemize}
È comprensibile che gli sviluppatori pensino solo ad arrivare al lancio e a
guadagnare abbastanza per continuare; ma ``la dura realtà della giungla legale''
colpisce dopo, quando si cerca di far rivivere il gioco.

\subsection{I tre enigmi}

Ogni volta che GOG cerca di riportare in vita un gioco deve risolvere
\textbf{tre enigmi}:
\begin{enumerate}
  \item \textbf{Legale.} Verificare tutto, capire perché il gioco è scomparso e
        se si può rimediare. Spesso è noioso (mail, siti, telefonate), ma a volte
        si arriva ad \textbf{assumere investigatori privati} per verificare che
        l'erede di una parte della licenza sia davvero chi si pensa; o a dover
        far riparlare due creatori ``dal carattere forte'' che non si parlano da
        vent'anni, perché decidano chi tiene i diritti. ``Di solito è un gran
        bussare alle porte.''
  \item \textbf{Economico.} GOG è un'azienda commerciale, parte di un gruppo
        quotato, che risponde a consiglio di amministrazione e azionisti: serve un
        \textbf{business case}. Anche quando tutti sono entusiasti, a volte non ha
        senso: nel caso di \emph{Black \& White}, per esempio, mettere al tavolo
        gli avvocati di Electronic Arts potrebbe costare più dei ricavi previsti
        in dieci anni\footnote{Il passaggio è confuso nella trascrizione (``black
        and white game, bringing the lawyers with electronic cards''); il senso
        riportato è quello ricostruibile.}.
  \item \textbf{Tecnico}, ``sorprendentemente il più facile''.
\end{enumerate}

\subsection{La preservazione giocabile}

GOG fa \textbf{preservazione giocabile}: il gioco non deve solo essere salvato
sul supporto originale, ma deve \textbf{girare sui PC moderni} e partire come se
lo aspetta un giocatore di oggi. Questo significa aggiungere funzioni di
\emph{quality of life}, a volte inesistenti all'epoca:
\begin{itemize}
  \item \textbf{salvataggi in cloud} (i giochi degli anni Novanta non avevano
        connessione a internet);
  \item supporto ai \textbf{doppi monitor} e alle \textbf{alte risoluzioni};
  \item supporto ai \textbf{gamepad}, ``tema caldo''.
\end{itemize}
GOG esiste da circa \textbf{18 anni} e ha sviluppato \textbf{strumenti interni}
che permettono di farlo su larga scala. Il risultato più importante è il
\textbf{GOG Preservation Program}: GOG si impegna a far funzionare i giochi del
programma sulle versioni moderne dei sistemi ``in ogni caso'', anche subentrando
al supporto dello sviluppatore e del publisher originali. Secondo Piotr è
probabilmente il più grande (o uno dei più grandi) cataloghi di giochi
\textbf{mantenuti attivamente} al mondo.

Nel \textbf{2024} GOG è entrata in EFGAMP per condividere con archivisti,
biblioteche, storici e amministrazioni nazionali ed europee le conoscenze
accumulate in 18 anni e la propria rete nell'industria: ``è tempo di
restituire''.

% ---------------------------------------------------------------------
\section{Perché preservare conviene anche all'industria}
\label{sec:c1-4-business}

La moderatrice chiede se alla base di tutto ci sia la \textbf{nostalgia}. Per
Piotr la nostalgia è \textbf{solo un fattore}: la preservazione è oggi una
\textbf{strategia di business}.
\begin{itemize}
  \item L'industria del videogioco è dove il cinema era decenni fa: i giochi sono
        investimenti enormi, e oggi si possono \textbf{rilanciare} dopo dieci o
        vent'anni, come si fa con i film.
  \item I giochi \textbf{invecchiano più lentamente} perché la tecnologia
        progredisce più lentamente: \emph{The Witcher 3} ha avuto una nuova
        edizione dopo dieci anni e ``sembra e si gioca ancora benissimo'';
        \emph{Oblivion} è stato rifatto dopo 18 anni (``se all'uscita eri troppo
        giovane, ora sei abbastanza grande per giocarci legalmente'').
  \item Tutti i grandi attori (PlayStation, Microsoft, Activision, Electronic
        Arts, Nintendo) hanno \textbf{unità di preservazione} che cercano e
        mettono al sicuro tutto ciò che trovano in azienda.
\end{itemize}

\paragraph{L'esempio della patch.} Si pensi a un gioco come \emph{The Last of
Us}: anni dopo esce la serie TV e la gente torna a giocare. Prima, per
\textbf{patchare} il gioco bisognava trovare qualcuno che si ricordasse come
fare, ricostruire in uno studio remoto un ambiente di lavoro, testare,
correggere, ritestare: circa \textbf{sei mesi}. Con tutto il materiale
preservato in un unico posto bastano \textbf{settimane}.

\paragraph{Asset per altri media.} I giochi generano altri prodotti culturali
(film, serie TV): ai creatori arrivano continuamente richieste del tipo ``avete
ancora questo asset del gioco, per usarlo nella serie?''.

\paragraph{La comunità come risorsa.} Se il gioco non è preservato e non è
giocabile, non sai come correggerlo e la sua \textbf{comunità muore}. ``Nell'era
dello sviluppo dell'IA, la comunità e la IP sono gli asset più forti''. Se il
gioco resta vivo puoi parlare con la comunità, capire se ha senso un
\textbf{remaster} o un \textbf{remake} e, se funziona, fare un
\textbf{sequel}, continuando a ottenere un ritorno sull'investimento.

% ---------------------------------------------------------------------
\section{Il futuro del supporto fisico}
\label{sec:c1-4-fisico}

Una spettatrice, Yasmin, pone ``la domanda di un grande sogno'': Sony avrebbe
annunciato che con la PlayStation 6 non ci saranno più dischi\footnote{La notizia
è riportata dalla spettatrice; non è verificata nel capitolo.}, anche per le
dimensioni dei giochi (cita le circa 260 GB di \emph{GTA VI}). Come si è passati
dal floppy alla chiavetta USB, all'HDD, all'SSD, potrebbe nascere un nuovo oggetto
fisico capace di contenere giochi così grandi e far tornare il mercato fisico,
anche solo per gli appassionati?

\paragraph{Willem, da archivista.} Il supporto fisico \textbf{non conta}: per la
preservazione l'oggetto digitale vale quanto quello fisico. La fine del disco era
inevitabile; i supporti, anche conservati in condizioni ideali nei depositi,
``diventano polvere'' in uno o due secoli, ``i migliori''. Il digitale è sempre
più facile da preservare a lungo.

\paragraph{Willem, da giocatore.} Non esclude che un supporto fisico resti, ma sarà
\textbf{di altissima fascia}, per ``collezionisti con la C maiuscola'': pagheresti
300 euro per una copia fisica di \emph{GTA VI}, solo il supporto, magari con la
scatola? Diventerà un \textbf{prodotto di lusso}, utile ai publisher anche per
tenere legata la comunità dei collezionisti. Quanto alle tecnologie:
\begin{itemize}
  \item il \textbf{nastro} non va sottovalutato: ha le capacità più alte,
        soprattutto con le nuove tecnologie;
  \item la memorizzazione su \textbf{DNA} è ancora sperimentale, ma esiste;
  \item sono soluzioni molto costose e specifiche: ``il 99\% dei consumatori non
        se ne curerà'' e passerà al digitale, anche se a molti presenti
        dispiacerà.
\end{itemize}

\paragraph{Piotr.} Su PC il mondo è digitale da almeno un decennio; su PlayStation
il passaggio è stato più lento, ma c'è stato; l'annuncio è solo ``un campanello
d'allarme''. Dal punto di vista dello sviluppo un ritorno è improbabile, perché il
digitale è più comodo ed economico:
\begin{itemize}
  \item non bisogna produrre e spedire dischi;
  \item non bisogna consegnare la \emph{gold master} tre mesi prima del lancio:
        si può lavorare fino all'ultimo secondo.
\end{itemize}
E già oggi, comprando il disco di un gioco PlayStation, la console ne copia il
contenuto sul disco interno: il disco è \textbf{solo una licenza}. Ciò che si
perde davvero è il \textbf{mercato dell'usato}.

GOG è orgogliosa delle sue regole di \textbf{proprietà digitale}: un gioco
comprato non scompare mai dalla libreria e si può scaricare un
\textbf{installer offline}. Dovrebbe essere lo standard: le istituzioni e le
organizzazioni come EFGAMP devono battersi per regole sulla \textbf{proprietà
digitale}, perché i consumatori devono avere la certezza che ciò che pagano non
sparirà. Un gioco può sparire dal negozio per ragioni di licenza, ma chi l'ha
comprato deve continuare a possederlo. ``Viviamo in tempi interessanti'', e con
l'IA le cose potrebbero cambiare rapidamente.

\paragraph{Andreas Lange: le periferiche.} Lange sposta lo sguardo dai supporti
alle \textbf{periferiche}. Sui supporti concorda con gli altri; ma ``siamo esseri
analogici'', e per giocare servono \textbf{interfacce fisiche}. Per preservare un
\textbf{gameplay autentico} vanno preservate anche quelle: emulare un gioco arcade
su un portatile è un'esperienza completamente diversa dal giocarlo sull'hardware
originale, con le sue interfacce dedicate di input e output. Oggi c'è ancora un
mercato di controller ``vecchio stile'' perché molti giocatori vogliono la
sensazione originale; in futuro dovranno essere le istituzioni a mantenerli
funzionanti.

% ---------------------------------------------------------------------
\section{I premi}
\label{sec:c1-4-premi}

Il panel si chiude con due premi del GAMM, che la moderatrice mostra in
videochiamata:
\begin{itemize}
  \item a \textbf{EFGAMP}, ``per essere la prima e unica federazione europea
        dedicata alla preservazione e alla promozione del videogioco come medium
        artistico e culturale, e per il suo impegno nella salvaguardia del
        patrimonio digitale interattivo europeo per le generazioni future'';
  \item ad \textbf{Andreas Lange}, ``per il suo ruolo pionieristico nella
        preservazione e nella valorizzazione museale delle opere interattive, e
        per il costante impegno istituzionale a livello europeo nel promuovere e
        riconoscere il videogioco come forma di espressione culturale''.
\end{itemize}
La moderatrice sottolinea che il confronto con la federazione è sempre utile al
museo per capire che cosa fare: ``dà uno scopo a tutto ciò che facciamo''.

% ---------------------------------------------------------------------
\section*{Punti chiave del capitolo}
\addcontentsline{toc}{section}{Punti chiave del capitolo}

\begin{itemize}
  \item EFGAMP riunisce 30 membri (dalle iniziative di base alle biblioteche
        nazionali): un unicum mondiale. Le \textbf{comunità} (emulatori) sono
        state decisive, ma servono \textbf{istituzioni} sostenibili e
        sistematiche.
  \item Problemi europei: \textbf{copyright frammentato} tra stati e passaggio
        dal supporto fisico ai \textbf{giochi su server}, che il detentore dei
        diritti può spegnere (campagna Stop Killing Games).
  \item \textbf{Games2Cloud} collega il patrimonio videoludico al Cultural
        Heritage Cloud europeo (oltre 102 milioni, 2023--2029); EFGAMP prototipa
        la prima \textbf{ontologia} per le collezioni di videogiochi.
  \item Segnali: \emph{Black Myth: Wukong} al padiglione cinese della Biennale;
        i videogiochi tra i tre settori di uno studio pilota europeo sul
        patrimonio nato digitale.
  \item Un archivio pubblico (Sound \& Vision) preserva in digitale, con il
        permesso dei detentori, e cerca anche il \textbf{contesto}; con
        \textbf{Control Save} offre linee guida agli sviluppatori perché
        preservino \textbf{prima} che conoscenze e materiali vadano perduti.
        Non esiste una regola d'oro su che cosa donare.
  \item GOG: per riportare in vita un gioco servono tre passi (\textbf{legale},
        \textbf{economico}, \textbf{tecnico}); i giochi scompaiono soprattutto
        per accordi persi e licenze a tempo. La \textbf{preservazione
        giocabile} aggiunge funzioni moderne; il GOG Preservation Program
        garantisce che i giochi funzionino sui sistemi attuali.
  \item Preservare conviene: rilanci, remaster, remake, sequel, patch in
        settimane invece che in mesi, asset per serie e film, comunità vive.
  \item Il supporto fisico diventerà un \textbf{lusso per collezionisti}; il vero
        nodo è la \textbf{proprietà digitale} (installer offline, nessuna
        sparizione dalla libreria). Ma le \textbf{periferiche} vanno preservate
        per un gameplay autentico.
\end{itemize}
